\chapter{Teoria aproksymacji}
\label{ma2:ch:approximation}

Rozwinięcie Taylora pozwala przybliżać funkcje
za pomocą wielomianów w otoczeniu ustalonego punktu.
Nie każda funkcja ciągła jest jednak różniczkowalna,
a nawet funkcja gładka nie musi być sumą swojego
szeregu Taylora. W tym rozdziale zmieniamy pytanie:
czy na całym przedziale domkniętym można przybliżyć
dowolną funkcję ciągłą wielomianem z dowolnie małym
błędem? Nie wymagamy już zgodności pochodnych
w wybranym punkcie.

\section{Aproksymacja jednostajna i moduł ciągłości}

\begin{definicja}[Błąd najlepszej aproksymacji]
\label{ma2:def:approx-best-error}
Niech $K\subseteq\R^d$ będzie niepustym zbiorem
zwartym, a $C(K)$ przestrzenią rzeczywistych funkcji
ciągłych z normą supremum
$\|f\|_\infty=\sup_{x\in K}|f(x)|$.
Dla niepustego zbioru $A\subseteq C(K)$ określamy
\[
 E(f,A)=\inf_{g\in A}\|f-g\|_\infty.
\]
Na przedziale $[a,b]$, $a<b$, oznaczamy przez
$\mathcal P_n$ przestrzeń ograniczeń wielomianów
stopnia co najwyżej $n$ do tego przedziału
i piszemy $E_n(f)=E(f,\mathcal P_n)$.
\end{definicja}

\begin{stwierdzenie}
\label{ma2:prop:approx-distance}
Dla dowolnego niepustego $A\subseteq C(K)$
zachodzi $|E(f,A)-E(h,A)|\le\|f-h\|_\infty$.
Ponadto $E(f,A)=0$ wtedy i tylko wtedy,
gdy $f$ należy do domknięcia $A$ w normie supremum.
Na przedziale ciąg $(E_n(f))_{n\ge0}$ jest
nierosnący.
\end{stwierdzenie}
\begin{proof}
Dla każdego $g\in A$ nierówność trójkąta daje
$\|f-g\|_\infty\le\|f-h\|_\infty+\|h-g\|_\infty$.
Biorąc infimum po $g$ i zamieniając role $f,h$,
otrzymujemy pierwszą tezę. Jeżeli $E(f,A)=0$,
to dla każdego $\varepsilon>0$ istnieje $g\in A$
o $\|f-g\|_\infty<\varepsilon$; jest to dokładnie
warunek przynależności do domknięcia.
Implikacja odwrotna wynika z tego samego warunku.
Ostatnia teza wynika z inkluzji
$\mathcal P_n\subseteq\mathcal P_{n+1}$.
\end{proof}

\begin{definicja}[Moduł ciągłości]
\label{ma2:def:approx-modulus}
Dla $f\in C(K)$ i $\delta\ge0$ określamy
\[
 \omega_f(\delta)=
 \sup\{|f(x)-f(y)|:x,y\in K,\ \|x-y\|_2\le\delta\}.
\]
\end{definicja}

\begin{stwierdzenie}
\label{ma2:prop:approx-modulus}
Moduł ciągłości jest niemalejący,
$\omega_f(0)=0$, $\omega_f(\delta)\le2\|f\|_\infty$
oraz $\omega_f(\delta)\to0$ przy $\delta\to0^+$.
Jeżeli $K$ jest wypukły, to dla $\delta>0$
i $x,y\in K$ zachodzi
\[
 |f(x)-f(y)|\le
 \left(1+\frac{\|x-y\|_2}{\delta}\right)\omega_f(\delta).
\]
\end{stwierdzenie}
\begin{proof}
Monotoniczność i wartość w zerze wynikają
z definicji, a górne ograniczenie z nierówności
trójkąta. Funkcja ciągła na zbiorze zwartym
jest jednostajnie ciągła, co daje granicę.
Jeśli $x\ne y$, dzielimy odcinek $[x,y]$
na $N=\lceil\|x-y\|_2/\delta\rceil$ równych
części. Wszystkie punkty podziału należą do $K$,
a kolejne odległości nie przekraczają $\delta$.
Sumowanie $N$ różnic i nierówność
$N\le1+\|x-y\|_2/\delta$ dają wynik.
Dla $x=y$ lewa strona jest zerowa.
\end{proof}

\section{Wielomiany Bernsteina}

\begin{definicja}
\label{ma2:def:approx-bernstein}
Dla $f\in C([0,1])$, $n\ge1$ i $0\le k\le n$
określamy
\[
 b_{n,k}(x)=\binom nk x^k(1-x)^{n-k},\qquad
 B_nf(x)=\sum_{k=0}^n f(k/n)b_{n,k}(x).
\]
$B_nf$ nazywamy wielomianem Bernsteina funkcji $f$.
\end{definicja}

Dla ustalonego $x\in[0,1]$ liczby $b_{n,k}(x)$
są nieujemne i sumują się do $1$.
Wartość $B_nf(x)$ jest więc średnią ważoną
wartości funkcji w węzłach $k/n$, a nie
na ogół wartością interpolującą w tych węzłach.
W punktach $0$ oraz $1$ odpowiednio jedynie
składnik $k=0$ i $k=n$ ma wagę $1$,
zatem $B_nf(0)=f(0)$ i $B_nf(1)=f(1)$.

\begin{lemat}[Momenty wag Bernsteina]
\label{ma2:lem:approx-bernstein-moments}
Dla $x\in[0,1]$ i $n\ge1$ mamy
\[
 \sum_{k=0}^n b_{n,k}(x)=1,\qquad
 \sum_{k=0}^n \frac{k}{n}b_{n,k}(x)=x,
\]
\[
 \sum_{k=0}^n
 \left(\frac{k}{n}-x\right)^2 b_{n,k}(x)
 =\frac{x(1-x)}{n}.
\]
\end{lemat}
\begin{proof}
Pierwsza tożsamość jest wzorem dwumianowym
dla $x+(1-x)$. Dla $k\ge1$ zachodzi
$k\binom nk=n\binom{n-1}{k-1}$, więc
$\sum_k k b_{n,k}(x)=nx$ po wyłączeniu $nx$
i ponownym zastosowaniu wzoru dwumianowego.
Dla $n\ge2$ i $k\ge2$ analogicznie
$k(k-1)\binom nk=n(n-1)\binom{n-2}{k-2}$,
stąd $\sum_k k(k-1)b_{n,k}(x)=n(n-1)x^2$.
Dla $n=1$ obie strony tej ostatniej równości
są zerowe. Korzystając z $k^2=k(k-1)+k$,
otrzymujemy
\[
 \sum_k (k/n)^2 b_{n,k}(x)
 =\frac{n-1}{n}x^2+\frac{x}{n}.
\]
Rozwinięcie kwadratu $(k/n-x)^2$ i podstawienie
poprzednich tożsamości daje tezę.
\end{proof}

\begin{twierdzenie}[Aproksymacja wielomianami Bernsteina]
\label{ma2:thm:approx-bernstein}
Dla $f\in C([0,1])$, $n\ge1$ i $\delta>0$
zachodzi
\[
 \|B_nf-f\|_\infty\le
 \left(1+\frac{1}{2\delta\sqrt n}\right)
 \omega_f(\delta).
\]
W szczególności
$\|B_nf-f\|_\infty\le\frac32\omega_f(n^{-1/2})$
i $B_nf$ zbiega jednostajnie do $f$.
\end{twierdzenie}
\begin{proof}
Na mocy \ref{ma2:lem:approx-bernstein-moments}
oraz nieujemności wag
\[
 |B_nf(x)-f(x)|\le
 \sum_{k=0}^n b_{n,k}(x)|f(k/n)-f(x)|.
\]
Stosując \ref{ma2:prop:approx-modulus},
otrzymujemy ograniczenie przez
\[
 \omega_f(\delta)\left(1+
 \frac1\delta\sum_{k=0}^n b_{n,k}(x)|k/n-x|\right).
\]
Nierówność Cauchy'ego--Schwarza dla skończonych
ciągów $\sqrt{b_{n,k}(x)}$ i
$\sqrt{b_{n,k}(x)}|k/n-x|$ oraz lemat o momentach
dają
\[
 \sum_k b_{n,k}(x)|k/n-x|
 \le\sqrt{\frac{x(1-x)}n}\le\frac1{2\sqrt n}.
\]
Ostatnie oszacowanie wynika z
$x(1-x)=1/4-(x-1/2)^2$.
Biorąc supremum po $x$ otrzymujemy pierwszą
nierówność. Podstawiamy $\delta=n^{-1/2}$
i stosujemy granicę modułu ciągłości z
\ref{ma2:prop:approx-modulus}.
\end{proof}

\begin{twierdzenie}[Twierdzenie Weierstrassa o aproksymacji]
\label{ma2:thm:approx-weierstrass}
Niech $a<b$. Dla każdej $f\in C([a,b])$
i każdego $\varepsilon>0$ istnieje wielomian $p$
taki, że $\|f-p\|_\infty<\varepsilon$.
Równoważnie $\lim_{n\to\infty}E_n(f)=0$.
\end{twierdzenie}
\begin{proof}
Funkcja $g(t)=f(a+(b-a)t)$ jest ciągła na $[0,1]$.
Z \ref{ma2:thm:approx-bernstein} wynika jednostajna
zbieżność $B_ng$ do $g$. Po podstawieniu
\[
 p_n(x)=(B_ng)\bigl((x-a)/(b-a)\bigr)
\]
otrzymujemy wielomiany stopnia co najwyżej $n$
oraz $\|p_n-f\|_\infty=\|B_ng-g\|_\infty$,
gdyż podstawienie afiniczne jest bijekcją
między przedziałami. Stąd $E_n(f)$ jest
ograniczony przez ciąg błędów dążący do zera.
Implikacja odwrotna wynika z definicji infimum.
\end{proof}

\begin{przyklad}[Aproksymacja funkcji kwadratowej]
\label{ma2:ex:approx-bernstein-square}
Dla $f(x)=x^2$ tożsamości momentów dają
\[
 B_nf(x)=x^2+\frac{x(1-x)}n,\qquad
 \|B_nf-f\|_\infty=\frac1{4n}.
\]
Maksimum błędu jest osiągane w $x=1/2$.
Dla $n\ge2$ funkcja $f$ sama należy do $\mathcal P_n$,
więc $E_n(f)=0$. Konstrukcja Bernsteina
nie musi zatem dawać wielomianu najlepszej
aproksymacji; jej zaletą jest jawny wzór
i gwarancja zbieżności dla każdej funkcji ciągłej.
\end{przyklad}

\begin{przyklad}[Funkcje lipschitzowskie]
\label{ma2:ex:approx-bernstein-lipschitz}
Jeżeli $|f(x)-f(y)|\le L|x-y|$ na $[0,1]$,
to bezpośrednio w pierwszym oszacowaniu dowodu
\ref{ma2:thm:approx-bernstein} dostajemy
\[
 |B_nf(x)-f(x)|\le
 L\sum_k b_{n,k}(x)|k/n-x|
 \le\frac{L}{2\sqrt n}.
\]
Przykładem jest $f(x)=|x-1/2|$ z $L=1$,
co pokazuje, że metoda działa także
dla funkcji nieróżniczkowalnych.
\end{przyklad}

\section{Wielomiany Bernsteina wielu zmiennych}

\begin{definicja}
\label{ma2:def:approx-bernstein-multi}
Niech $d\ge1$, $n\ge1$, $K=[0,1]^d$ oraz $f\in C(K)$.
Dla $k=(k_1,\ldots,k_d)\in\{0,\ldots,n\}^d$
i $x=(x_1,\ldots,x_d)\in K$ określamy
\[
 w_{n,k}(x)=\prod_{i=1}^d b_{n,k_i}(x_i),
 \qquad
 B_n^{(d)}f(x)=
 \sum_{k\in\{0,\ldots,n\}^d}f(k/n)w_{n,k}(x).
\]
Tutaj $k/n=(k_1/n,\ldots,k_d/n)$, a $b_{n,k_i}$
są wagami z definicji \ref{ma2:def:approx-bernstein}.
\end{definicja}

Jest to wielomian stopnia co najwyżej $n$
w każdej zmiennej z osobna, a zatem stopnia
całkowitego co najwyżej $dn$.
Wagi są nieujemne, a ich suma wynosi $1$:
przy rozpisaniu iloczynu skończonych sum
otrzymujemy
\[
 \sum_k w_{n,k}(x)
 =\prod_{i=1}^d\sum_{j=0}^n b_{n,j}(x_i)=1
\]
na mocy \ref{ma2:lem:approx-bernstein-moments}.
Wartość wielomianu jest więc średnią ważoną
wartości $f$ na regularnej siatce w sześcianie.

\begin{twierdzenie}
\label{ma2:thm:approx-bernstein-multi}
Dla $f\in C([0,1]^d)$, $n\ge1$ i $\delta>0$
zachodzi
\[
 \|B_n^{(d)}f-f\|_\infty\le
 \left(1+\frac{\sqrt d}{2\delta\sqrt n}\right)
 \omega_f(\delta).
\]
W szczególności
\[
 \|B_n^{(d)}f-f\|_\infty
 \le\frac32\omega_f\bigl(\sqrt{d/n}\bigr)
 \longrightarrow0
 \qquad(n\to\infty)
\]
dla ustalonego $d$.
\end{twierdzenie}
\begin{proof}
Dzięki faktoryzacji wag i
\ref{ma2:lem:approx-bernstein-moments}
mamy dla każdego $i\in\{1,\ldots,d\}$
\[
 \begin{aligned}
 \sum_k w_{n,k}(x)(k_i/n-x_i)^2
 &=
 \left(\sum_{j=0}^n b_{n,j}(x_i)(j/n-x_i)^2\right)
 \prod_{l\ne i}\sum_{j=0}^n b_{n,j}(x_l)\\
 &=\frac{x_i(1-x_i)}n.
 \end{aligned}
\]
Po zsumowaniu względem $i$ dostajemy
\[
 \sum_k w_{n,k}(x)\|k/n-x\|_2^2
 =\sum_{i=1}^d\frac{x_i(1-x_i)}n
 \le\frac d{4n}.
\]
Nierówność Cauchy'ego--Schwarza i równość
$\sum_k w_{n,k}(x)=1$ dają
\[
 \sum_k w_{n,k}(x)\|k/n-x\|_2
 \le\sqrt{\frac d{4n}}.
\]
Ponieważ sześcian jest wypukły, stosujemy
\ref{ma2:prop:approx-modulus}:
\[
 \begin{aligned}
 |B_n^{(d)}f(x)-f(x)|
 &\le\sum_k w_{n,k}(x)|f(k/n)-f(x)|\\
 &\le\omega_f(\delta)
 \sum_k w_{n,k}(x)
 \left(1+\frac{\|k/n-x\|_2}{\delta}\right)\\
 &\le\left(1+\frac{\sqrt d}{2\delta\sqrt n}\right)
 \omega_f(\delta).
 \end{aligned}
\]
Bierzemy supremum po $x$, a następnie
$\delta=\sqrt{d/n}$, co daje drugie oszacowanie.
Granica wynika z jednostajnej ciągłości $f$,
jak w \ref{ma2:prop:approx-modulus}.
\end{proof}

\begin{wniosek}
\label{ma2:cor:approx-bernstein-box}
Wielomiany są gęste w normie supremum
w przestrzeni funkcji rzeczywistych ciągłych
na każdym prostopadłościanie
$K=\prod_{i=1}^d[a_i,b_i]$, gdzie $a_i<b_i$.
\end{wniosek}
\begin{proof}
Odwzorowanie
$\Phi:[0,1]^d\to K$, określone przez
$\Phi_i(t)=a_i+(b_i-a_i)t_i$, jest bijekcją
afiniczną z odwrotnością o składowych
$(x_i-a_i)/(b_i-a_i)$.
Dla $f\in C(K)$ stosujemy
\ref{ma2:thm:approx-bernstein-multi} do $f\circ\Phi$.
Wielomiany
$p_n=(B_n^{(d)}(f\circ\Phi))\circ\Phi^{-1}$
są wielomianami w $x$ dzięki afiniczności
$\Phi^{-1}$. Bijektywność $\Phi$ zapewnia
$\|p_n-f\|_{\infty,K}
=\|B_n^{(d)}(f\circ\Phi)-f\circ\Phi\|_{\infty,[0,1]^d}$,
co kończy dowód.
\end{proof}

\begin{przyklad}[Aproksymacja powierzchni kwadratowej]
\label{ma2:ex:approx-bernstein-surface}
Dla $f(x,y)=x^2+y^2$ na $[0,1]^2$
faktoryzacja sum oraz momenty z
\ref{ma2:lem:approx-bernstein-moments} dają
\[
 B_n^{(2)}f(x,y)
 =x^2+y^2+\frac{x(1-x)+y(1-y)}n.
\]
Istotnie, część sumy z $(k_1/n)^2$
jest iloczynem drugiego momentu wag w $x$
i sumy wag w $y$, a część z $(k_2/n)^2$
obliczamy analogicznie.
Wobec tego błąd jest nieujemny i wynosi
\[
 \|B_n^{(2)}f-f\|_\infty=\frac1{2n},
\]
gdyż każda z dwóch funkcji $t(1-t)$
ma maksimum $1/4$ dla $t=1/2$.
Mimo że $f$ należy do przestrzeni wielomianów
już dla stopnia całkowitego $2$,
wielomiany Bernsteina nie odtwarzają jej
dokładnie dla żadnego skończonego $n$.
\end{przyklad}

\section{Twierdzenie Stone'a--Weierstrassa}

Poprzednie konstrukcje korzystają z regularnej siatki.
Chcemy teraz uzyskać twierdzenie o aproksymacji
także na zwartych zbiorach o dowolnym kształcie,
a ponadto przez rodziny funkcji inne niż wielomiany.
Wystarczy odpowiednia struktura algebraiczna
oraz możliwość rozróżnienia dwóch punktów.

\begin{definicja}
\label{ma2:def:approx-separate}
Niech $K$ będzie niepustą zwartą przestrzenią
Hausdorffa. Przez $C(K,\R)$ oznaczamy algebrę
ciągłych funkcji rzeczywistych na $K$ z działaniami
punktowymi i normą supremum.
Podalgebra $A\subseteq C(K,\R)$ rozdziela punkty,
jeżeli dla dowolnych różnych $x,y\in K$
istnieje $g\in A$ takie, że $g(x)\ne g(y)$.
Mówimy, że $A$ zawiera stałe, jeśli każda funkcja
stała należy do $A$.
\end{definicja}

\begin{lemat}
\label{ma2:lem:approx-lattice}
Niech $A\subseteq C(K,\R)$ będzie podalgebrą
zawierającą stałe, a $B$ jej domknięciem
w normie supremum. Wówczas $B$ jest podalgebrą,
zawiera stałe, a jeśli $u,v\in B$, to także
$|u|$, $\max(u,v)$ i $\min(u,v)$ należą do $B$.
\end{lemat}
\begin{proof}
Suma i mnożenie przez skalar są ciągłe w normie
supremum. Jeżeli $u_j,v_j\in A$ oraz
$u_j\to u$, $v_j\to v$ jednostajnie, to
\[
 \|u_jv_j-uv\|_\infty
 \le\|u_j\|_\infty\|v_j-v\|_\infty
    +\|v\|_\infty\|u_j-u\|_\infty.
\]
Ciąg $(\|u_j\|_\infty)$ jest ograniczony,
więc prawa strona dąży do zera.
Otrzymujemy $uv\in B$; pozostałe działania
algebraiczne zachowują $B$ przez ciągłość.
Stałe należą do $B$, bo należą do $A$.

Dla $u\in B$ wybierzmy $M>0$ tak, aby
$\|u\|_\infty\le M$. Na mocy
\ref{ma2:thm:approx-weierstrass} istnieją
wielomiany $p_j$ zbiegające jednostajnie do
$t\mapsto |t|$ na $[-M,M]$.
Ponieważ $B$ jest algebrą zawierającą stałe,
$p_j\circ u\in B$ dla każdego $j$.
Ponadto
\[
 \|p_j\circ u-|u|\|_\infty
 \le\sup_{|t|\le M}|p_j(t)-|t||\longrightarrow0.
\]
Domkniętość $B$ daje $|u|\in B$.
Dla liczb rzeczywistych $a,b$ zachodzi
$\max(a,b)=(a+b+|a-b|)/2$ i
$\min(a,b)=(a+b-|a-b|)/2$;
sprawdzamy to osobno dla $a\ge b$ i $a<b$.
Stosując te tożsamości punktowo, dostajemy
pozostałe dwie tezy.
\end{proof}

\begin{twierdzenie}[Twierdzenie Stone'a--Weierstrassa, wersja rzeczywista]
\label{ma2:thm:approx-stone-weierstrass}
Niech $K$ będzie niepustą zwartą przestrzenią
Hausdorffa, a $A\subseteq C(K,\R)$ podalgebrą
zawierającą stałe i rozdzielającą punkty.
Wówczas $A$ jest gęsta w $C(K,\R)$
w normie supremum.
\end{twierdzenie}
\begin{proof}
Niech $B$ będzie domknięciem $A$.
Ustalmy $f\in C(K,\R)$ i $\varepsilon>0$.

\emph{Krok 1. Dopasowanie w dwóch punktach.}
Dla $x,y\in K$ znajdziemy $h_{x,y}\in A$
takie, że $h_{x,y}(x)=f(x)$ i
$h_{x,y}(y)=f(y)$.
Dla $x=y$ wystarczy funkcja stała $f(x)$.
Jeśli $x\ne y$, wybieramy $g\in A$
z $g(x)\ne g(y)$ i kładziemy
\[
 h_{x,y}(z)=f(x)+
 \frac{f(y)-f(x)}{g(y)-g(x)}(g(z)-g(x)).
\]
Jest to kombinacja liniowa $g$ i funkcji stałej,
a podstawienie $z=x$ oraz $z=y$ daje
wymagane wartości.

\emph{Krok 2. Dolne oszacowanie globalne
przy zachowaniu wartości w punkcie $x$.}
Ustalmy $x\in K$. Dla każdego $y\in K$
zbiór
\[
 U_{x,y}=\{z\in K:h_{x,y}(z)>f(z)-\varepsilon\}
\]
jest otwarty dzięki ciągłości funkcji
i zawiera $y$. Zwartość daje skończone
podpokrycie $U_{x,y_1},\ldots,U_{x,y_r}$.
Definiujemy
\[
 g_x=\max_{1\le j\le r}h_{x,y_j}.
\]
Na mocy \ref{ma2:lem:approx-lattice}, zastosowanego
skończenie wiele razy, $g_x\in B$.
Każdy z maksymalizowanych składników
ma w $x$ wartość $f(x)$, więc $g_x(x)=f(x)$.
Dla dowolnego $z\in K$ istnieje $j$ z
$z\in U_{x,y_j}$, skąd
$g_x(z)\ge h_{x,y_j}(z)>f(z)-\varepsilon$.

\emph{Krok 3. Górne oszacowanie globalne.}
Dla każdego $x\in K$ zbiór
\[
 V_x=\{z\in K:g_x(z)<f(z)+\varepsilon\}
\]
jest otwarty i zawiera $x$.
Wybieramy skończone podpokrycie
$V_{x_1},\ldots,V_{x_s}$ i kładziemy
$h=\min_{1\le i\le s}g_{x_i}$.
Ponownie $h\in B$ dzięki
\ref{ma2:lem:approx-lattice}.
Wszystkie funkcje $g_{x_i}$ są większe
od $f-\varepsilon$ w każdym punkcie,
więc także $h>f-\varepsilon$.
Z pokrycia przez $V_{x_i}$ wynika
$h<f+\varepsilon$ w każdym punkcie.
Otrzymujemy $\|h-f\|_\infty\le\varepsilon$.

Ponieważ $\varepsilon>0$ było dowolne,
istnieje ciąg elementów $B$ zbiegający do $f$.
Domkniętość $B$ daje $f\in B$.
Stąd $B=C(K,\R)$.
\end{proof}

W dowodzie potrzebowaliśmy możliwości dopasowania
wartości w dwóch punktach, a następnie możliwości
tworzenia skończonych maksimów i minimów.
Pierwszą zapewnia rozdzielanie punktów wraz
z obecnością stałych. Drugą otrzymaliśmy
z własności algebraicznych i jednowymiarowego
twierdzenia Weierstrassa, przechodząc do domknięcia.

\begin{wniosek}[Aproksymacja na dowolnym zwartym podzbiorze przestrzeni euklidesowej]
\label{ma2:cor:approx-polynomials-compact}
Dla każdego niepustego zwartego $K\subseteq\R^d$,
$d\ge1$, ograniczenia rzeczywistych wielomianów
w $d$ zmiennych do $K$ są gęste w $C(K)$.
\end{wniosek}
\begin{proof}
Ograniczenia wielomianów tworzą podalgebrę
zawierającą stałe. Jeżeli $x,y\in K$
są różne, to dla pewnego $i$ zachodzi
$x_i\ne y_i$. Wielomian współrzędnościowy
$z\mapsto z_i$ rozdziela te punkty.
Stosujemy \ref{ma2:thm:approx-stone-weierstrass}.
\end{proof}

\begin{przyklad}[Aproksymacja na okręgu]
\label{ma2:ex:approx-circle}
Na okręgu $K=\{(x,y)\in\R^2:x^2+y^2=1\}$
możemy przybliżyć dowolną funkcję ciągłą
przez ograniczenia wielomianów w $x,y$
na mocy \ref{ma2:cor:approx-polynomials-compact}.
Na przykład funkcję $(x,y)\mapsto|y|$
możemy przybliżyć wielomianami w samym $y$,
stosując \ref{ma2:thm:approx-weierstrass}
do $t\mapsto|t|$ na $[-1,1]$.
Natomiast wielomiany w samej współrzędnej $x$
nie są gęste w $C(K)$: dla każdego takiego
wielomianu $p$ i funkcji $f(x,y)=y$
mamy
\[
 \|f-p\|_\infty\ge
 \max\{|1-p(0)|,|-1-p(0)|\}\ge1.
\]
Ostatnia nierówność wynika z nierówności
trójkąta, gdyż suma obu modułów jest
co najmniej $2$.
Wynika to z niemożności rozdzielenia
punktów $(0,1)$ i $(0,-1)$ przez tę podalgebrę.
\end{przyklad}

\begin{przyklad}[Rola obecności funkcji stałych]
\label{ma2:ex:approx-no-constants}
Podalgebra
$A=\{p\in\mathcal P:p(0)=0\}\subseteq C([0,1])$,
gdzie $\mathcal P$ oznacza wszystkie ograniczenia
wielomianów, rozdziela punkty, bo zawiera
funkcję $x\mapsto x$.
Nie jest jednak gęsta, gdyż dla każdego $p\in A$
mamy $\|1-p\|_\infty\ge |1-p(0)|=1$.
Nie można zatem w ogólnym twierdzeniu
pominąć warunku dotyczącego stałych bez
zastąpienia go odpowiednim dodatkowym założeniem.
\end{przyklad}
